\cvsection{Experiencia}

\cventrykeep
\cventry{05/2026 - Actual}{Ingeniero en IA}{Agrosuper}{consultor \textperiodcentered\ AI/ML Center of Excellence (CoE)}{agrosuper.png}{%
\begin{cvitems}
  \item Construí ML pipelines en Azure Databricks con Python, SQL y PySpark DataFrames, cubriendo development, training, validation, deployment, inference, monitoring y lifecycle management.
  \item Desarrollé jobs ETL y de feature processing distribuidos en Spark para datasets de entrenamiento y scoring, integrando validación de datos y Databricks Jobs programados.
  \item Trabajé dentro del AI/ML Center of Excellence (CoE), aportando a estándares de MLOps, guías de arquitectura, patrones ML reutilizables y estructuras estandarizadas de proyectos/pipelines.
  \item Definí y apliqué prácticas de validation, deployment y production readiness, incluyendo model monitoring, observability, reproducibility, traceability, versionado y documentación operativa.
  \item Impulsé secure development practices en Dev / QA / Prod: access control, secrets management, least privilege en credenciales, aislamiento de entornos y deployments controlados de activos ML.
  \item Acompañé a equipos en la adopción de prácticas consistentes de AI/ML engineering en Azure Databricks, combinando entrega hands-on de pipelines con governance y estandarización del model lifecycle.
\end{cvitems}%
}

\cventrykeep
\cventry{03/2026 - 05/2026}{Ingeniero en IA}{Accenture}{consultor}{accenture.png}{%
\begin{cvitems}
  \item Diseñé sistemas Agentic AI / Multi-Agent Systems sobre Azure y Databricks, con orquestadores que reciben inputs, invocan tools/APIs, llaman LLMs y devuelven structured outputs o acciones.
  \item Implementé pipelines GenAI en Python con Prompt Engineering, patrones LangChain/LangGraph, OpenAI APIs, tool calling y structured outputs para automatización de workflows.
  \item Construí flujos agenticos que combinan retrieval/tools/APIs con LLMs (RAG cuando aplica) para obtener contexto, ejecutar acciones y generar respuestas context-aware.
  \item Integré REST APIs y sistemas internos en orquestación de agentes: input $\rightarrow$ orchestrator/agent $\rightarrow$ retrieval/tools/APIs $\rightarrow$ LLM $\rightarrow$ structured response/action.
  \item Desplegué componentes de IA con CI/CD en Azure/Databricks, priorizando pipelines de agentes reproducibles frente al uso puntual de prompts contra un único LLM.
\end{cvitems}%
}

\cventrykeep
\cventry{01/2026 - 03/2026}{Científico de Datos}{Colmena}{consultor}{colmena.png}{%
\begin{cvitems}
  \item Migré 4 workloads ML/SQL productivos de Snowflake a Azure Databricks, adaptando lógica de negocio y ETL a notebooks Databricks y jobs Spark (lift-and-shift controlado).
  \item Reescribí SQL de Snowflake a PySpark/Spark SQL DataFrames: joins, aggregations, window functions (\texttt{ROW\_NUMBER}/\texttt{QUALIFY}), date transformations y upserts tipo \texttt{MERGE}.
  \item Traduje constructos Snowflake a equivalentes Spark (\texttt{IFF}/\texttt{DECODE} $\rightarrow$ \texttt{CASE}, filtros window, lógica de fechas) preservando paridad funcional de los modelos migrados.
  \item Ejecuté validaciones de migración: row counts, null checks, detección de duplicados, reconciliación de métricas de negocio y chequeos funcionales antes del cutover.
  \item Optimicé ETL y workloads ML en Azure Databricks con Python, SQL, PySpark y Spark SQL para production readiness post-migración.
  \item Preparé notebooks/pipelines migrados para ejecución productiva en Azure, validando salidas ML y readiness operativa de los cuatro workloads.
\end{cvitems}%
}

\cventrykeep
\cventry{05/2025 - 12/2025}{Ingeniero en IA y Automatización}{Unilever}{consultor}{unilever.png}{%
\begin{cvitems}
  \item Construí pipelines de document processing en Python para múltiples tipos de documentos (PDF y otras fuentes no estructuradas): ingestion, text extraction, parsing y estructuración de campos.
  \item Implementé documents $\rightarrow$ ingestion $\rightarrow$ extraction $\rightarrow$ parsing $\rightarrow$ structuring $\rightarrow$ normalization $\rightarrow$ transformation $\rightarrow$ report generation.
  \item Automatizé normalización y transformación con expresiones regulares y salidas tabulares/tipo Excel para alimentar reporting operativo.
  \item Integré extracción y reporting vía REST APIs y JSON sobre Azure, orquestando la entrega automatizada de reportes a partir de datos documentales estructurados.
  \item Fortalecí el pipeline de automatización documental con validación de campos extraídos y ejecuciones reproducibles sobre lotes de documentos en Azure.
\end{cvitems}%
}

\cventrykeep
\cventry{01/2025 - 04/2025}{Desarrollador en Automatización con IA}{Megatronix}{consultor}{megatronix.png}{%
\begin{cvitems}
  \item Construí agentes conversacionales y automatizaciones en n8n como plataforma central de orquestación para interacciones de leads/soporte en WhatsApp e Instagram.
  \item Diseñé flujos channel event $\rightarrow$ webhook $\rightarrow$ n8n workflow $\rightarrow$ HTTP Request nodes / REST APIs $\rightarrow$ LLM $\rightarrow$ business logic $\rightarrow$ response/action.
  \item Orquesté procesos de leads y customer support con triggers n8n, webhooks, payloads JSON e integraciones LLM API con Prompt Engineering.
  \item Conecté canales de mensajería a sistemas downstream mediante nodos HTTP/API de n8n, enrutando eventos hacia respuestas automatizadas y acciones de seguimiento.
  \item Entregué automatización centrada en nodos n8n, webhooks, REST APIs y agentes conversacionales con LLM.
\end{cvitems}%
}

\cventrykeep
\cventry{02/2024 - 12/2024}{Científico de Datos}{Banco Galicia}{consultor}{Banco_Galicia.png}{%
\begin{cvitems}
  \item Desarrollé modelos de merchant risk y fraud detection en Azure con Python, Pandas, NumPy y Scikit-learn (Random Forest), usando SMOTE para classification desbalanceada.
  \item Ingesté y procesé datos de Brighterion y Cybersource: consolidación de variables de riesgo, cleaning y Feature Engineering para scoring de comercios.
  \item Entrené y validé modelos de classification con threshold analysis para balancear detección frente a falsos positivos en fraud scoring.
  \item Combiné scores ML con hard rules / business rules en sistemas de decisión para merchant risk y control de fraude por contracargos.
  \item Implementé fraud monitoring y métricas productivas en Azure (evaluación orientada a precision/recall y seguimiento operativo).
\end{cvitems}%
}

\cventrykeep
\cventry{08/2023 - 01/2024}{Científico de Datos}{AB InBev}{consultor}{logo_AB_InBev.png}{%
\begin{cvitems}
  \item Construí pipelines distribuidos de feature generation y procesamiento de audiencias con PySpark DataFrames y Spark SQL en Azure Databricks (joins, aggregations, ETL reutilizable).
  \item Generé datasets customer-level para marketing mediante cleaning, transformaciones Spark y feature generation orientadas a segmentación y audiencias personalizadas.
  \item Automatizé notebooks/jobs de Databricks desde la ingesta hasta outputs listos para campañas de marketing.
  \item Analicé performance de segmentación/audiencias y campaign analytics, presentando resultados y recomendaciones a stakeholders.
  \item Entregué analytics orientada a Data Engineering en Azure Databricks con Python, SQL y Pandas para preparación y evaluación de datasets.
\end{cvitems}%
}

\cventrykeep
\cventry{04/2021 - 07/2023}{Científico de Datos}{Mercado Libre}{}{MELI.jpeg}{%
\begin{cvitems}
  \item Lideré un equipo de 3 personas de Data Science en prevención de fraude, diseñando pipelines de risk scoring con Python, SQL, Pandas, NumPy y Scikit-learn sobre AWS.
  \item Utilicé Amazon S3 para almacenamiento de datasets, AWS Glue para ETL, Athena para SQL analítico, EMR para procesamiento distribuido y SageMaker para experimentación, training y deployment de ML.
  \item Implementé supervised classification y Feature Engineering para fraud/risk scoring, combinando salidas ML con reglas de fraude en sistemas de decisión de alto volumen.
  \item Desarrollé fraud profiling y segmentación con PCA, K-Means y HDBSCAN para identificar patrones de comportamiento y clusters en datos transaccionales.
  \item Construí detección de image similarity con OpenCV, SIFT/SURF y feature matching para análisis de patrones en imágenes asociadas a fraude.
  \item Operé scoring y monitoring en SageMaker/Jupyter para transacciones de alto volumen: validation, scoring productivo y seguimiento de métricas de fraude/riesgo.
\end{cvitems}%
}

\cventrykeep
\cventry{08/2019 - 03/2021}{Científico de Datos}{Ixpandit}{}{ixpandit.png}{%
\begin{cvitems}
  \item Desarrollé scorecards de Credit Risk con Logistic Regression / supervised classification (Python, Pandas, NumPy, Scikit-learn) para clientes nuevos y renovadores.
  \item Ejecuté data $\rightarrow$ EDA/cleaning $\rightarrow$ Feature Engineering $\rightarrow$ training $\rightarrow$ validation $\rightarrow$ scorecard/scoring $\rightarrow$ monitoring, incluyendo tratamiento de missing values y selección de variables.
  \item Construí ETL y pipelines de scoring en SQL/Python para decisiones de crédito, con validación de modelos y monitoreo de métricas de classification en producción.
  \item Contenericé componentes de scoring/validación con Docker y realicé experimentación en Jupyter Notebook.
  \item Monitoreé performance de scorecards en el tiempo y recalibré pipelines de scoring para carteras de nuevos vs.\ renovadores ante drift de métricas.
\end{cvitems}%
}

\cventrylegacy{09/2018 - 07/2019}{Analista de Control de Gestión}{Inworx}{}{inworx.jpg}{%
\begin{cvitemslegacy}
  \item Brindé soporte a proyectos, control de gestión y reportes para áreas de negocio.
\end{cvitemslegacy}%
}

\cventrylegacy{10/2015 - 08/2018}{Analista de Gestión de Riesgos}{QBE Seguros La Buenos Aires}{}{qbe.jpeg}{%
\begin{cvitemslegacy}
  \item Evalué riesgos corporativos, diseñé planes de mitigación y di seguimiento a su ejecución.
\end{cvitemslegacy}%
}

\cventrylegacy{03/2013 - 09/2015}{Analista Administrativo}{Rosa Ban Plásticos}{}{rbp.jpg}{%
\begin{cvitemslegacy}
  \item Apoyé reportes operativos y planificación de flujos administrativos.
\end{cvitemslegacy}%
}
\par\vspace{0.6em}%
